\section{The Kazakh-FEVER 3K Benchmark}
\label{sec:benchmark}

\subsection{Corpus Construction and Annotation Protocol}
To construct \textbf{Kazakh-FEVER 3K}, we established a curated knowledge repository comprising over 250,000 passage chunks extracted from the Kazakh Wikipedia dump and contemporary public fact-checking repositories (including Factcheck.kz). Each text chunk was filtered for syntactic completeness, normalized for Cyrillic-Kazakh orthographic consistency, and indexed for rapid retrieval.

Annotation was conducted by a team of native Kazakh linguists and professional fact-checkers under a rigorous three-stage protocol:
\begin{enumerate}[leftmargin=*,nosep]
    \item \textbf{Claim Generation}: Annotators extracted factual assertions from verified reference passages. For each reference passage, annotators generated claims belonging to three mutually exclusive classes: \textsc{Supported} (claim is fully entailed by evidence), \textsc{Refuted} (claim is explicitly contradicted), and \textsc{Not Enough Info} (claim contains unverified factual components).
    \item \textbf{Evidence Alignment}: Annotators manually highlighted minimal sentences in candidate corpora required to substantiate or disprove claims. Evidential dependencies were recorded with sentence-level granularity.
    \item \textbf{Quality Control and Agreement}: All samples underwent cross-validation by independent secondary annotators. Inter-annotator agreement on a randomly sampled 20\% subset yielded Cohen's Kappa $\kappa = 0.84$ and Fleiss' Kappa $0.81$, reflecting substantial annotator alignment.
\end{enumerate}

\subsection{The Hard NEI Protocol}
A pervasive weakness of automated verification datasets is model reliance on lexical overlap heuristics for the \textsc{NEI} class. In standard FEVER benchmarks, NEI instances are frequently created by pairing claims with unrelated passages, leading models to learn that low vocabulary overlap reliably signals NEI.

To eliminate this spurious shortcut, we formulated the \textbf{Hard NEI} protocol. Annotators construct NEI claims maintaining over 85\% lexical and structural overlap with genuine candidate passages, but modify subtle semantic operators that cannot be confirmed from the text:
\begin{itemize}[leftmargin=*,nosep]
    \item \textbf{Evidentiality and Epistemic Modality}: Altering direct past suffixes (\textit{-ды/-ді}) to inferential or quotative forms (\textit{-ыпты/-іпті}, \textit{екен}) or modal markers (\textit{мүмкін}, \textit{тиіс}).
    \item \textbf{Quantifier and Scope Alternations}: Shifting universal quantifiers (\textit{барлық}, \textit{барша}) to existential or partial quantifiers (\textit{кейбір}, \textit{кейде}).
    \item \textbf{Temporal and Causal Anchoring}: Introducing unstated dates, temporal adverbs, or causal subordinators (\textit{-ғандықтан}) that are grammatically plausible but ungrounded in evidence.
\end{itemize}
Consequently, models cannot determine veracity via lexical similarity alone, but must analyze whether the precise semantic proposition is entailed by the retrieved passage.

\begin{table}[t]
\centering
\small
\caption{Distribution of claim-evidence pairs across splits in the Kazakh-FEVER benchmark.}
\label{tab:dataset_stats}
\resizebox{\columnwidth}{!}{%
\begin{tabular}{lrrrr}
\toprule
\textbf{Split} & \textbf{Supported} & \textbf{Refuted} & \textbf{Hard NEI} & \textbf{Total} \\
\midrule
Train & 274 & 154 & 210 & 638 \\
Development & 58 & 33 & 44 & 135 \\
Test & 61 & 33 & 47 & 141 \\
\midrule
\textbf{Overall} & \textbf{393} & \textbf{220} & \textbf{301} & \textbf{914} \\
\bottomrule
\end{tabular}%
}
\end{table}

The dataset statistics are summarized in Table~\ref{tab:dataset_stats}. While the broad benchmark design targets 3,000 balanced pairs (2,100 training, 450 development, and 450 test pairs split into 300 encyclopedic and 150 emerging instances), our verified release establishes 914 rigorously checked claim-evidence pairs across 71 articles in eight domains (\textit{history}, \textit{geography}, \textit{literature}, \textit{law}, \textit{society}, \textit{arts}, \textit{technology}, \textit{science}), partitioned into 638 train (69.8\%), 135 dev (14.8\%), and 141 evaluation pairs (15.4\%). An auxiliary cross-domain benchmark of 1,000 synthetic social media posts across five platforms and eight personas enables testing out-of-distribution robustness.
